\section{Metrics}
\label{appendix:metrics}

We employ different evaluation protocols for the two task categories:

For \textbf{Deep Reasoning Tasks}, we follow standard QA evaluation practices as ToG \citep{sun2023tog,ma2024tog2} and HippoRAG \citep{hipporag, gutiérrez2025HippoRAG2}:
\begin{itemize}
    \item \textbf{Exact Match (EM)}: Measures the percentage of predictions that exactly match the ground truth answer. Specifically, we follows the Substring-based EM metric (used in ToG/ToG-2~\citep{sun2023tog,ma2024tog2}) to robustly assess answer accuracy in longer, natural-language response generated by LLMs, which goes through the whole response to check whether the answer is in. 
    \item \textbf{F1 Score}: Computes word-level overlap between predictions and ground truth answers.
\end{itemize}

For \textbf{Broad Reasoning Tasks}, we adopt a multi-dimensional LLM-based evaluation approach due to the complexity and open-ended nature of these queries following LightRAG \citep{guo2024lightrag}:
\begin{itemize}
    \item \textbf{Comprehensiveness (Comp.)}: Measures how thoroughly the answer addresses all aspects of the question.
    \item \textbf{Diversity (Div.)}: Assesses the variety of perspectives and insights provided in the answer.
    \item \textbf{Empowerment (Emp.)}: Evaluates how well the answer enables informed understanding and judgment.
\end{itemize}
The LLM-based evaluation uses GPT-4o-mini as judge, with careful attention to prompt design and answer ordering to avoid positional bias. The LLM evaluation prompt is shown in Appendix.\ref{appendix:prompts}